\section{Conclusions}
\label{sec:conclusions}

This paper has presented a metric-consistent Riemannian pseudospectral decomposition framework for tracking structural regimes within the nocturnal stable boundary layer. By utilizing a fractional hyperbolic coordinate stretching function under an economy-size thin SVD architecture, we project sparse vertical tower data onto a continuous computational manifold. This structure isolates scale-aware diagnostic signatures without introducing ad-hoc cutoff frequencies or non-local Gibbs ringing.

Applying this architecture to the CASES-99 campaign under a GMM segmentation schema reveals three statistically separable operational regimes, supported by a silhouette coefficient peak at $K=3$: Continuous Turbulence, Wave-Dominated Stable, and Intermittent Shear Bursts. We further frame these diagnostics in a reduced-order interpretation motivated by fast/slow SBL time-scale separation ($\epsilon = \tau_t/\tau_m \ll 1$), map the informational metrics to Ozmidov-scale constraints (Eq.~\ref{eq:d_eff_ozmidov_scaling}) as a consistency relation, formulate a ``Virtual Tower'' interpolation operator ($\mathcal{P}_{\mathrm{VT}}$) for NWP model verification, and construct a $2/3$-rule spectral anti-aliasing filter (Eq.~\ref{eq:exponential_filter}) to ensure numerical stability in forward implementations.

This approach shifts stable boundary layer analysis away from localized gradient parameters and toward geometry-aware manifold diagnostics, offering a new path for robust sub-mesoscale parameterization and climate model closure development. The strong suppression of unphysical upper-boundary wave reflections ($\mathcal{R} \sim 3.7 \times 10^{-6}$) indicates that this pipeline can be exported to other dense micrometeorological tower domains (e.g., SHEBA, FLUXNET) with appropriate revalidation. This enables physically interpretable, automated regime classification with direct utility for data assimilation, turbulence closure development, and high-fidelity mesoscale model validation.

% ============================================================================
